\section{Keeping the operations when we change the object}
\label{m:section}

The common comparison has produced intrinsic simple factors. Its group example
also showed a limit: the factors of $C_4$ and $C_2\times C_2$ agree, although
the groups do not. We now develop that question in a setting where the
attachments can be calculated explicitly. The objects will be modules,
and our main examples will be vector spaces equipped with an operator.
We retain the operator in every map, just as the dimension argument
retained scalar multiplication in every quotient.

\subsection{Modules and their maps}

A ring $R$ has an abelian group under addition, an associative multiplication
with identity $1$, and the two distributive laws. Multiplication need not be
commutative. A left $R$-module is an abelian group $M$ with operations $rx$
for $r\in R$, $x\in M$, satisfying
\[
 r(x+y)=rx+ry,\quad (r+s)x=rx+sx,\quad
 r(sx)=(rs)x,\quad 1x=x.
\]
The same cancellation arguments as for vector spaces give $r0=0$,
$0x=0$ and $(-1)x=-x$. A submodule contains zero and is closed under
addition and all these scalar operations. A module map preserves addition
and satisfies $f(rx)=rf(x)$. In this part all unspecified maps between
modules are module maps. We work over a fixed ring whenever forming
extensions or comparing modules.

Intersections of submodules are submodules. Their sum $U+W$ consists of
the elements $u+w$; distributivity proves that it is a submodule, and it
is the smallest one containing both. The direct sum of two modules is
the set of pairs with coordinatewise addition and scalar action. Each
module identity holds in both coordinates. Its inclusions and projections
are module maps. A map out of $A\oplus B$ is determined by its restrictions
to $(a,0)$ and $(0,b)$, since $(a,b)=(a,0)+(0,b)$.

For a submodule $N\subseteq M$, the additive quotient $M/N$ is already
defined by the group construction. Set $r(x+N)=rx+N$. Replacing $x$ by
$x+n$ changes $rx$ by $rn\in N$, so this action is well-defined;
the module identities descend on representatives. The projection is a
module map. If a module map $f:M\to X$ kills $N$, its unique additive
factorization $\bar f:M/N\to X$ is a module map, because
\[
 \bar f(r(x+N))=f(rx)=rf(x)=r\bar f(x+N).
\]
Thus the quotient universal property retains every operation required here.

Kernels, images and preimages of submodules are submodules: apply a map
to a sum or a scalar multiple, and use the corresponding closure in the
target. An image of a submodule is the image of the map restricted to it.
Consequently the first isomorphism theorem supplies a module isomorphism
\[
 M/\Ker f\iso\im f.
\]
The inverse of a bijective module map is a module map: writing $y=f(x)$
gives $f^{-1}(ry)=rx$, and the analogous calculation handles sums.
All the short exact sequences below mean injectivity, surjectivity, and
image equal to kernel for these maps.

\subsection{Preimages, quotients, and the local comparison}

For a surjection $q:M\to Q$ with kernel $N$, inverse image and direct
image give an inclusion-preserving correspondence
between submodules of $Q$ and submodules of $M$ containing $N$.
Indeed, for a submodule $H$ of $M$,
\[
 q^{-1}(q(H))=N+H:
 \quad q(x)=q(h)\ \Longleftrightarrow\ x-h\in N.
\]
Every quotient map and induced isomorphism in this correspondence is
linear over $R$. For $N\subseteq L\subseteq M$ this also proves
$(M/N)/(L/N)\cong M/L$ by $(x+N)+(L/N)\mapsto x+L$.

The kernel calculation at the heart of the butterfly remains valid:
if $A_1\subseteq A_2$ and $B_1\subseteq B_2$ are submodules of $M$,
write $C=A_2\cap B_2$, $U=A_1\cap B_2$, $W=A_2\cap B_1$.
The map
\[
 C\longrightarrow (A_1+C)/(A_1+W),\qquad c\longmapsto c+(A_1+W)
\]
is onto, and its kernel is $U+W$. To check the nontrivial inclusion,
if $c=a+w$ with $c,w\in C$ and $a\in A_1$, then $a=c-w\in U$.
The reverse inclusion follows by substitution. Exchanging $A$ and $B$
gives the same kernel and the other wing. This is a concrete calculation
of the common support quotient, retaining the $R$-action. It is available independently
of any theorem about composition series.

\subsection{Finite-dimensional tools, after dimension has been proved}
\label{m:finite-tools}

The later operator examples use subspace bases and lifts. We establish
their existence here, after Theorem~\ref{e:dimension-theorem}; none is
being fed back into the earlier proof of dimension.

\begin{lemma}\label{m:bases}
If a vector space $V$ has a finite basis of size $n$, every subspace has
a finite basis. Every independent list in $V$ extends to a basis of $V$.
For a subspace $W$, a basis of $W$ followed by lifts of a basis of $V/W$
is a basis of $V$.
\end{lemma}
\begin{proof}
Start with any finite independent list. Pass through a fixed finite
spanning list of $V$, appending a vector exactly when it is outside
the span of the vectors already present. Each append preserves
independence: a relation with nonzero coefficient on the new vector
would put it in the old span. At the end the list spans $V$, so it is
a basis. The already proved invariance of basis size shows that the
original independent list had at most $n$ members.
An arbitrary independent set cannot contain $n+1$ distinct vectors,
so it too is finite and the same conclusion applies.

To construct a basis of $W$, start with the empty list and, as long
as it does not span $W$, choose an element outside its span. This
cannot produce $n+1$ independent vectors. Thus it stops after finitely
many steps with a basis of $W$. This requires only finitely many
choices, bounded by $n$.

The images in $V/W$ of a finite basis of $V$ span the quotient.
Applying the append procedure starting with the empty list to those
images gives a quotient basis. Choose one lift of each of its finitely
many members. The lifts together with a basis of $W$ span $V$:
subtract the lifted expression of a vector's quotient class, leaving
an element of $W$. In a relation among this list, projection to $V/W$
first forces all lift coefficients to vanish; independence in $W$
then forces the remaining coefficients to vanish.
\end{proof}

In particular, $\dim V=\dim W+\dim(V/W)$. Applying the first
isomorphism theorem to a linear map gives
$\dim V=\dim\Ker f+\dim\im f$. These are consequences of the
dimension theorem already proved, and will be used only from here on.
The same basis-lifting proof gives a vector-space section of any
surjective linear map onto a finite-dimensional space. Such a section
need not preserve a further module or operator structure.

\subsection{An operator is a polynomial action}
\label{m:operators}

For a field $k$, $k[t]$ consists of finite formal sums
$\sum_{i\geq0}a_it^i$. Addition is coefficientwise and multiplication
is convolution: the coefficient of $t^r$ in a product is
$\sum_{i+j=r}a_ib_j$. All sums involved are finite. Associativity
holds because either parenthesization of a triple product has coefficient
$\sum_{i+j+\ell=r}a_ib_jc_\ell$; distributivity is the corresponding
distributivity of finite sums. The constant polynomial $1$ is the identity.

Given a $k$-linear endomorphism $T$ of $V$, define
\[
 p(t)v=p(T)v=\sum_i a_iT^iv,\qquad T^0=\id.
\]
Induction gives $T^iT^j=T^{i+j}$, so convolution proves
$p(T)q(T)=(pq)(T)$. The other module axioms follow from linearity.
Conversely, a $k[t]$-module has its underlying $k$-vector space and
the linear operator $T(v)=tv$; the module axioms recover the displayed
polynomial action. These constructions are inverse.

A linear map $f:(V,T)\to(V',T')$ is a module map exactly when
$fT=T'f$. Necessity uses the polynomial $t$; sufficiency follows by
induction on powers and addition. Submodules are exactly the
$T$-stable subspaces. Thus the square
\[
\begin{tikzcd}[column sep=large,row sep=large]
V\arrow[r,"T"]\arrow[d,"f"']&V\arrow[d,"f"]\\
V'\arrow[r,"T'"']&V'
\end{tikzcd}
\]
records the extra condition that distinguishes an operator-preserving
map from an arbitrary linear map.

For $\lambda\in k$, let $S_\lambda$ be $k$ with $t$ acting by
$\lambda$. It is simple as a module because its only $k$-subspaces
are zero and itself. It is also $k[t]/(t-\lambda)$: evaluation
$p\mapsto p(\lambda)$ is onto, and the identity
$t^i-\lambda^i=(t-\lambda)\sum_{j=0}^{i-1}t^{i-1-j}\lambda^j$
shows that its kernel is $(t-\lambda)$.

We also use $P_r=k[t]/(t^r)$ for $r\geq1$. Its vectors have unique
representatives of degree less than $r$, by separating the terms of
degree at least $r$. Thus $1,t,\ldots,t^{r-1}$ is a vector-space
basis, and multiplication by $t$ shifts along a nilpotent chain.
No polynomial factorization theorem is required for these modules.

\subsection{Simple layers and finite length}
\label{m:length-foundations}

A module is simple if it is nonzero and has only zero and itself as
submodules. A composition series is a finite strictly increasing
submodule filtration from zero to the whole module with simple
successive quotients. A module has finite length if it admits such
a series; at this stage that phrase is an existence condition.
The zero module has the series with one term and no factors.

\begin{lemma}\label{m:simple-map}
A nonzero module map between simple modules is an isomorphism.
\end{lemma}
\begin{proof}
Its image is nonzero, hence is the whole target. Its kernel cannot be
the whole source, hence is zero. The kernel test gives injectivity,
and the inverse of a bijective module map preserves the module operations.
\end{proof}

Here a $k$-algebra is a ring $R$ with a unital map from $k$ into its
center. Restricting scalars makes each $R$-module a $k$-vector space,
and every $R$-submodule is a $k$-subspace.
Every finite-dimensional module over a $k$-algebra has a composition
series. If it is nonzero, choose a proper submodule of largest possible
vector-space dimension; zero is one candidate. The possible dimensions
form a finite set by Lemma~\ref{m:bases}. Its quotient is simple:
an intermediate nonzero proper submodule would lift to a larger proper
submodule. Repeat inside the chosen proper submodule. Dimensions
strictly decrease, so the process ends. This proves existence without
using uniqueness of composition factors.

We will also need closure of finite-length modules under submodules
and quotients over a general ring. Let
$0=M_0\subsetneq\cdots\subsetneq M_n=M$ be a composition series
and $N\subseteq M$ a submodule. Set
$N_i=N\cap M_i$ and $Q_i=(M_i+N)/N$.
For each $i$, inclusion and projection give a short exact sequence
\begin{equation}\label{m:layer-exact}
0\longrightarrow N_i/N_{i-1}\longrightarrow M_i/M_{i-1}
\longrightarrow Q_i/Q_{i-1}\longrightarrow0.
\end{equation}
The first map sends $n+N_{i-1}$ to $n+M_{i-1}$ and is injective.
The second sends $m+M_{i-1}$ to $(m+N)+Q_{i-1}$ and is onto.
Its kernel consists of $m=a+n$ with $a\in M_{i-1}$ and $n\in N$;
then $n=m-a\in N_i$, proving equality with the first image.
The simple middle module forces precisely one outer module to be zero
and the other to be isomorphic to it. Deleting repetitions from the
$N_i$ and $Q_i$ therefore gives composition series for $N$ and $M/N$.

Conversely, if $A$ and $B$ have composition series and
$0\to A\to E\xrightarrow{q}B\to0$ is short exact, identify $A$
with its image in $E$. Append to its series the preimages under $q$
of a composition series of $B$. The quotient correspondence proves
strictness, and the maps
\[
 q^{-1}(B_j)/q^{-1}(B_{j-1})\longrightarrow B_j/B_{j-1},\qquad
 e+q^{-1}(B_{j-1})\longmapsto q(e)+B_{j-1}
\]
are isomorphisms by their kernels and surjectivity. This constructs
a composition series for $E$. Its factors are exactly the two given
lists together. No uniqueness theorem is used in this construction.

We now have the objects and operations needed for an independent
spectral comparison of two module filtrations. Length and multiplicities
are already intrinsic by Theorem~\ref{j:jordan-holder}. The following
page calculations will compare the factors directly, without using that
theorem as a premise.
